% =========================
% puzzles/pXX.tex
% Final - The Talon Plates
% =========================

\PuzzleChapterIntro{The Talon Plates}{Tiling \textbullet\ Local Forcing \textbullet\ One Vault Code}{Do not be deceived by the area. Just because the numbers fit the ledger does not mean the plates fit the floor. You must find the 'atom' of the tiling to separate the true chambers from the impossible ones.}

\Lore{
The Archive Casino keeps a repair corridor: a brass runway scored into unit squares.

The steward slides a tray of identical talon plates---six unit-squares fused into a crooked claw.
At the bend of every claw is a pocket: a deliberate void the metal leans around but does not cover.

Some chambers accept the claws cleanly; others reject them no matter how patiently the pieces are turned.
The copper strip below does not ask \textit{how}; it asks only which chambers can be sealed.
}

\Problem
A \textbf{talon plate} is the fixed six-square shape below (filled squares are covered; the $\circ$ is the \emph{pocket square} and is \emph{not} covered):

\[
\begin{array}{ccc}
\blacksquare & \blacksquare & \blacksquare\\
\blacksquare & \circ         & \blacksquare\\
\blacksquare & \phantom{\blacksquare} & \phantom{\blacksquare}
\end{array}
\]

A talon may be translated, rotated, or reflected.

A rectangle of size $m\times n$ is called \textbf{sealable} if its $mn$ unit squares can be tiled by talon plates with no overlaps and no gaps.

\medskip
The copper strip lists the following chambers:

\[
\begin{array}{c|c}
\text{Rectangle} & \text{Rune}\\\hline
14\times18 & 11\\
8\times15 & 17\\
10\times30 & 19\\
12\times7 & 29\\
6\times20 & 37\\
22\times18 & 31\\
16\times9 & 41\\
12\times11 & 43\\
18\times18 & 23
\end{array}
\]

\VaultDirective{Sum the runes of the sealable chambers.}

\SolutionPage
\scriptsize
We give two necessary conditions (local forcing and a mod-$4$ obstruction), then tile exactly the chambers that pass.

\textbf{Linked-pair involution $\Rightarrow 12\mid mn$.}

Fix a tiling of an $m\times n$ rectangle.
For any plate $H$, let $P(H)$ be its \emph{pocket square}: the unique unit square adjacent to three covered squares of $H$.
Let $M(H)$ be the \emph{mouth square} of that pocket: the unique side-adjacent neighbor of $P(H)$ not covered by $H$.

Because the rectangle is perfectly tiled, $P(H)$ is covered by exactly one other plate; call it $f(H)$.
(So $f(H)\ne H$.)

\noindent\emph{Local forcing:}
If $K=f(H)$ covers $P(H)$, then $K$ must also cover $M(H)$ (it is the only available neighbor of $P(H)$ not already blocked by $H$).
Checking the finitely many placements of a talon that cover two adjacent squares $P(H)$ and $M(H)$ without overlapping $H$,
one finds that the only placements extendable to a full tiling are exactly those in which $H$ covers the pocket square $P(K)$.
Equivalently, $f(f(H))=H$.

Thus the plates partition into \emph{linked pairs} $\{H,f(H)\}$.
Each linked pair covers $12$ squares, so $mn$ must be a multiple of $12$.

\textbf{If $12\mid mn$ but $4\nmid m$ and $4\nmid n$, then the rectangle is not sealable.}

Assume $12\mid mn$ and $4\nmid m,4\nmid n$.
Then $m\equiv n\equiv 2\pmod 4$, so the number of linked pairs is
\[\frac{mn}{12},\]
which is odd.

A direct finite check shows that, up to rotation/reflection, the union of any linked pair is one of the following four $12$-square shapes:
\[
\begin{array}{c@{\qquad}c@{\qquad}c@{\qquad}c}
\text{(I)}\ \begin{array}{cccc}
\blacksquare&\blacksquare&\blacksquare&\blacksquare\\
\blacksquare&\blacksquare&\blacksquare&\blacksquare\\
\blacksquare&\blacksquare&\blacksquare&\blacksquare
\end{array}
&
\text{(II)}\ \begin{array}{ccc}
\blacksquare&\blacksquare&\blacksquare\\
\blacksquare&\blacksquare&\blacksquare\\
\blacksquare&\blacksquare&\blacksquare\\
\blacksquare&\blacksquare&\blacksquare
\end{array}
&
\text{(III)}\ \begin{array}{cccc}
\phantom{\blacksquare}&\blacksquare&\blacksquare&\phantom{\blacksquare}\\
\blacksquare&\blacksquare&\blacksquare&\blacksquare\\
\blacksquare&\blacksquare&\blacksquare&\blacksquare\\
\phantom{\blacksquare}&\blacksquare&\blacksquare&\phantom{\blacksquare}
\end{array}
&
\text{(IV)}\ \begin{array}{cccc}
\phantom{\blacksquare}&\blacksquare&\blacksquare&\blacksquare\\
\phantom{\blacksquare}&\blacksquare&\blacksquare&\blacksquare\\
\blacksquare&\blacksquare&\blacksquare&\phantom{\blacksquare}\\
\blacksquare&\blacksquare&\blacksquare&\phantom{\blacksquare}
\end{array}
\end{array}
\]

\noindent Now color the rectangle by making every $4$th \emph{column} black (columns $4,8,12,\dots$).
Since $n\equiv 2\pmod4$, the number of black columns is $\lfloor n/4\rfloor$, so the total number of black squares is
\[m\cdot\lfloor n/4\rfloor,\]
and since $m$ is even, this number is even.

In a linked pair of type (I) or (III), the shape spans $4$ consecutive columns and has exactly $3$ squares in each of those columns,
so it covers exactly $3$ squares in the unique black column among them---an \emph{odd} number.
In a linked pair of type (II) or (IV), every column contains an even number of squares, so it covers an \emph{even} number of black squares.

Therefore the number of linked pairs of type (I)/(III) must be even.
But the total number of linked pairs is odd, so the number of linked pairs of type (II)/(IV) is odd.

Now repeat the argument with every $4$th \emph{row} colored black.
The total number of black squares is $n\cdot\lfloor m/4\rfloor$, again even.
Pairs of type (II)/(IV) span $4$ consecutive rows and have $3$ squares in each row, so each such pair covers an odd number of black squares,
while pairs of type (I)/(III) cover an even number.
Hence the number of type (II)/(IV) pairs must be even---contradiction.

So if $12\mid mn$, at least one of $m,n$ must be divisible by $4$.

\textbf{Constructions.}

From type (I) above, two talons tile a $3\times4$ rectangle.
Therefore any rectangle whose side lengths are multiples of $3$ and $4$ is sealable by partitioning it into $3\times4$ blocks.

Also, since $12$ is a multiple of $3$ and $4$, any $12\times t$ rectangle with $t$ a sum of $3$'s and $4$'s is sealable by splitting into $12\times3$ and $12\times4$ strips.
In particular:
\[
12\times7=(12\times3)+(12\times4),\qquad
12\times11=(12\times3)+(12\times4)+(12\times4).
\]

\textbf{Evaluate the nine chambers.}

All nine rectangles have area divisible by $12$.
By Step 2, any rectangle with both side lengths $\equiv 2\pmod4$ is not sealable.
Thus
\[
14\times18,\ 10\times30,\ 22\times18,\ 18\times18\quad\text{are not sealable.}
\]

The remaining five are sealable by Step 3:
\[
8\times15,\ 12\times7,\ 6\times20,\ 16\times9\ \text{(multiples of $3$ and $4$)},\quad\text{and}\quad 12\times11.
\]

So the required rune-sum is
\[
17+29+37+41+43=167.
\]

\FinalAnswer{167}

\NextPage